% frontmatter/resumen.tex
% Resumen en castellano — máximo 1 página.
% ============================================================================
%
% GUÍA: El resumen debe cubrir:
%   1. Contexto del problema (corrección de exámenes en educación superior)
%   2. Qué se ha desarrollado (backend del SCDEE)
%   3. Tecnologías principales (Django, OCR, sistema distribuido)
%   4. Resultados/contribución principal
%   5. Conclusión breve
%
% Extensión recomendada: 150-250 palabras.
% ============================================================================

% TODO: Redactar el resumen. Ejemplo de estructura:
%
% La corrección de exámenes escritos en centros universitarios con gran
% volumen de alumnado es un proceso manual, lento y propenso a errores...
%
% Este trabajo presenta el diseño e implementación del backend de un
% Sistema de Corrección Digital de Exámenes Escritos (SCDEE), que permite
% digitalizar, procesar mediante OCR, distribuir y corregir exámenes de
% forma colaborativa...
%
% El sistema se ha desarrollado como una arquitectura distribuida basada en
% Django y Django REST Framework, con colas de tareas asíncronas (Celery),
% almacenamiento de objetos (MinIO), reconocimiento óptico de caracteres y
% una base de datos PostgreSQL...
%
% Los resultados demuestran que...

[Resumen pendiente de redacción.]

\palabrasclave{corrección de exámenes, OCR, sistema distribuido, Django, procesamiento de imágenes, API REST, backend}
